% ============================================================
% front.tex —— 前置部分：内容简介、前言、符号表
% 底本：LNM 2386
%   · 内容简介：PDF p2
%   · 前言（含致谢）：PDF p3–p6（印刷 v–viii）
%   · 符号表：PDF p15–p18（印刷 xvii–xx）
% 不译 Contents（PDF p7–13）
% ============================================================

% [原文 p.2]
\chapter*{内容简介}
\addcontentsline{toc}{chapter}{内容简介}

本书从统一的视角全面介绍谱网络，架起几何与超对称规范理论物理之间的桥梁。它提供了进入这一迅速发展领域前沿所需的基础背景，并平行处理几何与物理两个方面。在综述代数与几何中的基本主题之后，书中展开对高阶泰希米勒理论的详细介绍，包括用于 Hitchin 表示的 Fock–Goncharov 理论、极大表示以及较新的 $\Theta$ 正性概念。随后引入谱网络，强调它们借助所定义的非阿贝尔化与阿贝尔化映射在研究特征簇方面的效用。与此同时，书中探讨具有八个超荷的四维规范动力学的要点，包括电磁对偶、Seiberg–Witten 理论与 S 类理论。接着考察谱网络作为确定与分析 S 类理论中 BPS 谱之框架的作用。最后一章概述谱网络在一系列当代研究领域中的近期应用。本书面向希望进入该领域的研究人员与高年级学生，无论其背景是数学还是物理。

% [Tutor Note: 原文「eight supercharges」（八个超荷）即四维 $\mathcal N=2$ 超对称（每个超对称荷为四分量 Majorana 旋量，共 8 个超荷）；本书正文一律用花体 $\mathcal N$ 表示超对称的数目，勿与符号表中花体 $\mathcal N$ 的「量子场论中超对称的数目」混为一谈。]

% [原文 p.3（印刷 p.v）]
\chapter*{前言}
\addcontentsline{toc}{chapter}{前言}

谱网络由 Davide Gaiotto、Gregory W. Moore 与 Andrew Neitzke 在同名文章（2013 年发表于《\textit{Annales Henri Poincaré}》）中引入。它们是曲面上轨道的网络，自然出现在各种四维 $\mathcal N=2$ 量子场论的研究中，并提供一种特别有效的计算 BPS 谱的方法。从纯几何的观点看，它们给出黎曼曲面上的平坦联络与该曲面的一个谱覆盖上的平坦阿贝尔联络之间的一个新映射。同时，这些轨道网络给出平坦联络模空间上的局部坐标系，这些坐标系对于研究 Vladimir V. Fock 与 Alexander B. Goncharov 意义下的高阶泰希米勒空间很有价值。

自引入以来，谱网络不断证明自己是几何与数学物理中的一个深刻工具，其应用、推广以及与其他概念的联系与日俱增。然而，据我们所知，书本文献中仍缺少对该主题标准的、既严谨又具阐述性的综述；我们认为这主要是因为研究它们需要相当多样的背景。受我们各自多样研究兴趣的驱动，我们力求写出一本书：它既汇集了受过训练的数理学家与物理学家为严格入门该主题所需的必要基础，同时又展示了大量几何与数学物理中的应用，并以当前研究为指向。

本书旨在作为该主题的入门教材，以统一的方式从几何与数学物理的视角呈现谱网络及其众多应用。我们并不试图把这两种进路分开处理，而是提供宽广的动机与必要的背景，以抵达这一快速发展研究领域的前沿，同时解释相关的几何与物理方面。目标读者是希望进入该领域、具有数学或物理背景的研究人员与高年级学生。我们力求使这些几何与物理方面之间的联系更加透明，方法是在本书各节中有机地连接各个主题。鉴于谱网络的跨学科性质，我们尽力把书中所讨论的各个主题尽可能地统一起来。这种做法也特别解释了我们在同一卷书中同时使用标准数学术语与物理术语的决定。

% [原文 p.4（印刷 p.vi）]
谱网络及其在几何与数学物理中的应用是一个蓬勃发展的现代主题，当前正经历迅速的发展与扩张。因此，这一领域在工作中的数学家或物理学家看来，其面貌可能仍会发生变化。鉴于这种情况，我们选择以清晰的方式呈现理论的基本方面与迄今最坚实的应用。同时，我们收集了一系列当前研究发展的较广泛主题，在书的末尾以较为不严谨的方式呈现。这些方向上的重大进展可能使相关的进阶概念在未来以某种形式经历修订。

在正文之前的导论部分，我们提供一个历史概述，回顾高能物理、几何、特别是泰希米勒理论中的关键思想——正是它们导致谱网络作为计算具有两个或八个超对称的四维规范理论中 BPS 谱的工具被引入，以及它们在高等泰希米勒空间背景下的应用。虽然这部分并非作为这些主题的教学式导引而写——已有许多优秀的导引——但我们力求使尚不熟悉该领域的读者也能读懂并乐在其中，从而人人都能吸收关键概念，并从这一阐述中受益以理解本书后续各章。此处出现的一些重要数学概念会在后面各节中得到严谨而有教学意义的定义；我们在可能时强调这些联系。第 I 部分与第 II 部分完全独立于这一导论性论述，专注于几何中的各种主题。谱网络及其在几何中的用途在第 III 部分引入，而在第 IV 部分则从超对称规范理论的角度引入。

进一步展开本书的内容：第 I 部分是讨论谱网络所需数学背景的严谨而平实的入门。这里所包含的并非所有概念都在学术数学或物理课程中得到系统而充分的探讨，然而关于相关主题已有多种标准阐述；因此这一部分旨在为读者（无论其背景如何）建立共同的基础。对于已经熟悉所涉主题及其语言的读者，这一部分可以安全跳过。

第 II 部分是本书的核心部分之一，严谨地描述本书所讨论的高阶泰希米勒理论的多种主题。特别地，我们通过 Olivier Guichard 与 Anna Wienhard 引入的 $\Theta$ 正性这一核心概念，把高阶泰希米勒理论作为经典泰希米勒理论的高阶推广给出统一处理。此外，我们一方面讨论与泰希米勒空间上的 Thurston 剪切坐标的类比，另一方面讨论高阶泰希米勒空间上的 Fock–Goncharov 坐标。

% [原文 p.5（印刷 p.vii）]
谱网络在本书第 III 部分被明确定义。随后我们呈现并讨论谱网络在几何中最突出的应用：非阿贝尔化与阿贝尔化，它们把基础曲面的高阶泰希米勒空间与某个分歧覆盖——谱曲线——的阿贝尔特征簇联系起来。阿贝尔化的一种更解析的进路由精确 WKB 方法给出。从这个角度看，谱网络表现为 Borel 求和的 Stokes 图。Fock–Goncharov 坐标可以（至少在低秩时）由显式的单值数据恢复。

第 IV 部分探讨谱网络在具有 $\mathcal N=2$ 超对称的四维规范理论、特别是 S 类理论中的作用，将其作为描述 BPS 谱的有力框架。我们首先回顾四维超对称规范理论的一般方面，然后把注意力收窄到具有 $\mathcal N=2$ 超对称的理论。接着讨论电磁对偶并介绍 Seiberg–Witten 理论。最后引入 S 类理论并考察这些理论中 BPS 态的计数，特别是通过使用谱网络来进行。本部分意在作为该主题的一般性介绍，旨在为读者配备参与相应文献所需的基础概念。

第 V 部分呈现谱网络在几何与超对称量子场论两方面更高级的应用。与前面各章相比，这里的讨论更为非正式，应更多看作对进一步研究这些主题的邀请。我们提供了一份广泛的参考文献列表，鼓励读者查阅以获取更详细的论述。

为了尽量减少从事谱网络的数学界与数学物理界之间的语言障碍，本书以一份术语表作结，它对齐并澄清两个学科中所使用的核心概念。

本书源于作者们就此所涉主题组成的一个工作组，其起因是 2020 年 10 月 5 日至 9 日在法国 Luminy 的 CIRM 举办的会议「Teichmüller Theory: Classical, Higher, Super and Quantum」。在那次会议期间我们很快意识到：当时书本文献中缺少对谱网络的教学式而全面的介绍，并且该研究领域分裂为两个主要群体，二者之间的交流因语言与文化差异而受阻，尽管研究对象本质上是相同的。我们各自不同的研究背景在随后工作组的互动中被证明特别富有成效，使我们得以拥抱谱网络的各个方面，最终产生了把工作材料提炼为一本面向不同既有视角、旨在架起桥梁的实用手册的想法。

作者们感谢 CIRM 会议的组织者与参与者们的有益讨论，也感谢 CIRM 本身以及马赛 Calanques 的宁静而富有启发的自然风光促成了这个非常有益的工作组。

% [Tutor Note: CIRM 是 Centre International de Rencontres Mathématiques（国际数学会议中心）的缩写，位于法国马赛的 Luminy 校区；Calanques 是马赛沿海的石灰岩峡湾群。原文未展开缩写，译文照录并在此补注，方便读者查证会议与地点。]

% [原文 p.6（印刷 p.viii）]
我们感谢斯特拉斯堡大学与海德堡大学以及波恩的马克斯·普朗克数学研究所提供了肥沃的环境，本书的大部分工作在其中完成；感谢亚历山大·冯·洪堡基金会持续的支持，使我们得以在两所研究所之间开展多次研究访问。V.T. 感谢 R. Argurio、S. Banerjee、M. Bertolini、O. Hulik、L. Mol、A. Pasternak、S. Rota、R. Senghaas、R. Vandepopeliere 与 J. Walcher 的富有启发的讨论以及对本文早期版本的反馈。我们也感谢所有审稿人对我们工作的仔细审阅以及提出的一系列建设性意见与更正。

C.K. 获得了欧盟「地平线 2020」研究与创新计划下欧洲研究理事会（ERC）的资助（资助协议编号 101018839）。G.K. 感谢帕特雷大学科学委员会通过「Medicos」计划给予的支持。E.R. 由斯特拉斯堡大学的 Labex IRMIA 资助，并得到德国学术交流中心（DAAD）的博士后奖学金、欧盟「地平线 2020」研究与创新计划下欧洲研究理事会（ERC）的资助（资助协议编号 101018839），以及韩国高等研究院（KIAS）的个人资助（MG100901）。V.T. 感谢德国研究基金会（DFG）在德国卓越战略 EXC 2181/1 - 390900948（海德堡 STRUCTURES 卓越集群）下的资助。A.T. 的工作部分由 ERC-Advanced Grant 101018839、德国研究基金会（DFG）项目编号 281071066 - TRR 191 资助，并在项目末期由克洛德·贝尔纳里昂第一大学资助。

\begin{flushright}
海德堡与斯特拉斯堡，\\[2pt]
2023 年 7 月
\end{flushright}

% [原文 p.15（印刷 p.xvii）]
\chapter*{符号表}
\addcontentsline{toc}{chapter}{符号表}

\small
\begin{longtable}{@{}p{0.27\linewidth}p{0.67\linewidth}@{}}
$G$ & Lie 群 \\
$\g$ & Lie 代数 \\
$G^{\vee}$ & Langlands 对偶群 \\
$[G,G]$ & 换位子群 \\
$G_{ab}$ & 群的阿贝尔化 \\
$\mathrm{ad}$ & 伴随表示 \\
$\h$ & 复李代数的 Cartan 子代数 \\
$H$ & Cartan 子群 \\
$\g=\mathfrak{k}\oplus\mathfrak{p}$ & 实李代数的 Cartan 分解 \\
$\mathfrak{a}$ 或 $\mathfrak{t}$ & 实李代数的 Cartan 子代数 \\
$W$ & Weyl 群 \\
$K$ & Killing 形式 \\
$BG$ & 李群 $G$ 的分类空间 \\
$\g_\alpha$ & 相应于一个根的根空间 \\
$\operatorname{rad}\g$ & 李代数的根基 \\
$B$ 与 $\mathfrak{b}$ & Borel 子群及其李代数 \\
$U$ 与 $\mathfrak{u}$ & 上幂幺子群及其李代数 \\
$P$ 与 $\mathfrak{p}$ & 抛物子群及其李代数 \\
$\Lambda^{\perp}$ & 特征格 \\
$\Lambda$ & 余特征格 \\
$\mathcal{R}(G,T)$ & 根数据 \\
$\mathcal{R}(G,T)^{\vee}$ & 对偶根数据 \\
$\Sigma^{+}$ & 正根集 \\
$\Delta$ & 单根集 \\
$\Theta\subset\Delta$ & 抛物子群对应的单根子集 \\
$P_\Theta$ & 由单正根集 $\Theta$ 决定的抛物子群 \\
$\Sigma(\g,\alpha)$ & 限制根系 \\
% [原文 p.16（印刷 p.xviii）]
$S$ & 拓扑连通定向曲面 \\
$S_{g,h}$ & 亏格 $g$、$h$ 个边界分支的拓扑曲面 \\
$S_{g,\mathbf{p}}$ & 亏格 $g$、纤毛向量 $\mathbf{p}$ 的纤毛曲面 \\
$X$ & 常指黎曼曲面 \\
$\Sigma$ & 谱曲线 \\
$\tilde S$ & $S$ 的万有覆叠 \\
$\chi(S)$ & 曲面的欧拉示性数 \\
$P$ & 主丛 \\
$d_A$ & 向量丛或主丛上的联络 \\
$A$ & 联络矩阵 \\
$F(A)$ & 联络的曲率 \\
$P_{c(t)}$ & 沿曲线 $c(t)$ 的平行移动 \\
$\mathcal{L}$ & 局部系统 \\
$\mathcal{M}$ & 模空间 \\
$\mathcal{E}(S,G)$ & 基本群表示模去共轭作用所得的空间 \\
$\mathcal{R}(S,G)$ & 特征簇 \\
$\mathcal{R}^{\mathrm{fr}}(S,G)$ & 带标架表示的特征簇 \\
$\mathcal{R}^{\mathrm{fr}}_+(S,G)$ & 正带标架表示的特征簇 \\
$\mathcal{T}(S)$ & 泰希米勒空间 \\
$\mathcal{T}_H(S,G)$ & $G$-Hitchin 分支 \\
$\mathcal{X}_{S,G}$ & 带标架局部系统的模空间 \\
$\mathcal{A}_{S,G}$ & 装饰扭局部系统的模空间 \\
$\mathcal{M}^{\mathrm{fr}}(S,G)$ & 极大带标架表示的模空间 \\
$K_X$ & 复流形 $X$ 的典范线丛 \\
$\omega_{MC}$ & Maurer–Cartan 形式 \\
$\HH^2$ 或 $\HH$ & 双曲平面 \\
$\partial_\infty\HH$ & 双曲平面的无穷远边界 \\
$\Gamma$ & 图 \\
$\gamma$ & 曲面上的曲线 \\
$\mathcal{P}(S)$ & 曲面上的道路空间 \\
$\mathbb{Z}^{\mathrm{tw}}_{B'}[\mathcal{P}(\Sigma)]$ & 扭路径代数 \\
$\mathcal{W}_\vartheta$ & 角为 $\vartheta$ 的 WKB 谱网络 \\
$[A]$ & 某对象 $A$ 的等价类 \\
$X//G$ & GIT 商 \\
$\{\cdot,\cdot\}$ & 泊松括号 \\
$\mathcal{L}[f]$ & 函数的 Laplace 变换 \\
$S[f]$ & Borel 可和函数的 Borel 和 \\
$\operatorname{Res}_z\omega$ & 微分在一点处的留数 \\
$\widehat{\mathcal{F}}_\omega$ & 微分 $\omega$ 的叶状结构 \\
$\boldsymbol{\Phi}^{*h}$ & Hermite 共轭 \\
$\mathbb{T}$ & 标准三角形 \\
$T$ & 欧几里得三角形 \\
$\Delta$ & 曲面的三角剖分 \\
$E(\Delta)$ & 三角剖分的边 \\
$F(\Delta)$ & 三角剖分的面 \\
% [原文 p.17（印刷 p.xix）]
$E_e(\Delta)$ & 三角剖分的外边 \\
$E_i(\Delta)$ & 三角剖分的内边 \\
$V(\Delta)$ & 三角剖分的顶点 \\
$\mathrm{Lag}$ & 拉格朗日格拉斯曼流形 \\
$[L_1,M,L_2]$ & 三个两两横截的拉格朗日子流形构成的 Maslov 形式 \\
$T_\rho$ & 表示 $\rho$ 的 Toledo 数 \\
$\mathrm{Mod}(S_{g,\mathbf{p}})$ & 纤毛曲面的映射类群 \\
$\mathrm{PMod}(S_{g,\mathbf{p}})$ & 纤毛曲面的纯映射类群 \\
$\mathcal{F}$ & （部分）旗的空间 \\
$Q_\Delta$ & 与三角剖分 $\Delta$ 相关联的箭图 \\
$\mathcal{S}$ & 种子 \\
$\mathcal{X}_\Lambda$ & 种子的簇环面 \\
$\mu_i(Q)$ & 箭图 $Q$ 在顶点 $i$ 处的变异 \\
$\mathcal{A}_{\mathcal{S}}$ & 种子 $\mathcal{S}$ 的 A-簇代数簇 \\
$\mathcal{X}_{\mathcal{S}}$ & 种子 $\mathcal{S}$ 的 X-簇代数簇 \\
$\mathcal{N}$ & 量子场论中超对称的数目 \\
$\mathcal{H}$ & Hilbert 空间 \\
$\langle\cdot\,|\,\cdot\rangle$ & Hilbert 空间上的 Hermite 内积 \\
$|\phi\rangle$ & Hilbert 空间中的向量 \\
$\langle\phi|$ & 与向量相伴的线性泛函 \\
$\langle\mathcal{O}(x)\rangle$ & 局域场 $\mathcal{O}(x)$ 的关联函数 \\
$\mathcal{Z}$ & 中心荷 \\
$\mathcal{Q}$ & 超荷 \\
$\mathcal{Q}^I_\alpha$ & 超对称的生成元 \\
$P_\mu$ & 时空平移的生成元 \\
$M_{\mu\nu}$ & 时空旋转的生成元 \\
$\delta_{ij}$ & Kronecker 函数 \\
$\varepsilon_{ij}$ & 反对称符号（Levi–Civita 符号） \\
$\gamma^\mu$ & Dirac gamma 矩阵 \\
$\psi$ & 旋量 \\
$\sigma_i$ & Pauli 矩阵 \\
$[\cdot,\cdot]$ & 换位子 \\
$\{\cdot,\cdot\}$ & 反换位子 \\
$a$ 与 $a^{\dagger}$ & 湮灭与产生算子 \\
$\star$ & Hodge 星算子 \\
$S$ & 作用量泛函 \\
$\mathcal{L}$ & 拉格朗日量 \\
$E$ & 电场 \\
$B$ & 磁场 \\
$\rho_m$ & 磁荷密度 \\
$\mathbf{J}_m$ & 磁流 \\
$q_{\mathrm{el}}$ & 电荷 \\
$q_{\mathrm{mag}}$ & 磁荷 \\
% [原文 p.18（印刷 p.xx）]
$g$ & 非阿贝尔耦合常数 \\
$C(R)$ & 表示 $R$ 的 Casimir 数 \\
$I(R)$ & 表示 $R$ 的 Dynkin 指数 \\
$S_{\mathrm{YM}}(A)$ & 联络的 Yang–Mills 泛函 \\
$\mathcal{T}_g(\mathcal{M}_{(d-1)})$ & 与子流形 $\mathcal{M}_{(d-1)}$ 相关联的拓扑算子或缺陷 \\
$\mathcal{C}$ & 库仑分支 \\
$\lambda$ & Seiberg–Witten 微分（Liouville 典范 1-形式） \\
$\mathcal{S}[\mathcal{C},\g,D]$ & S 类理论 \\
$\mu$ & BPS 指标 \\
$F(\wp,\vartheta)$ & 带标架 $2d$–$4d$ 指标的生成函数 \\
$\mathcal{F}_G(X,\theta)$ & 指数化的 Chern–Simons 不变量 \\
$\mathrm{Vir}_c$ & Virasoro 代数 \\
$\mathrm{Heis}$ & Heisenberg 代数 \\
$\Psi$ & 共形块 \\
$m_Z$ & $Z$ 玻色子的能量标度 \\
$\Lambda_{\mathrm{QCD}}$ & 量子色动力学的动力学标度 \\
$\mathcal{L}_{\wp,\vartheta}$ & 超对称界面 \\
$\Psi_q$ & 量子双对数 \\
BPS & Bogomol'nyi–Prasad–Sommerfield \\
CFT & 共形场论（conformal field theory） \\
QCD & 量子色动力学（quantum chromodynamics） \\
QED & 量子电动力学（quantum electrodynamics） \\
QFT & 量子场论（quantum field theory） \\
SCFT & 超共形场论（superconformal field theory） \\
TQFT & 拓扑量子场论（topological quantum field theory） \\
VEV & 真空期望值（vacuum expectation value） \\
WKB & Wentzel–Kramers–Brioullin \\
\end{longtable}
\normalsize

% [Tutor Note: 符号表中 $G^{\vee}$ 的上标用的是逻辑「或」符号 $\vee$（Langlands 对偶），不是字母 v；$\g=\mathfrak{k}\oplus\mathfrak{p}$ 中 $\mathfrak{k}$ 是紧部分、$\mathfrak{p}$ 是非紧部分（Cartan 分解），原文花体与小写字母靠字体区分，排版时务请保留。]

% [Tutor Note: 原文把限制根系写作 $\Sigma(\g,\alpha)$（第二个变元是希腊字母 $\alpha$），而文献中更常见的是 $\Sigma(\g,\mathfrak{a})$（$\mathfrak{a}$ 为实李代数的极大分裂 Cartan 子代数）。此处照录原文，供读者对照文献时留意。]

% [Tutor Note: 符号 $\tilde S$ 在符号表中表示曲面 $S$ 的万有覆叠，而在书末术语表「Spectral network」条目中又用来表示谱覆盖 $\tilde S\to S$；二者同形异义，须按上下文区分（第 13–16 章中两种用法都出现）。]

% [Tutor Note: 符号表中 $g$ 出现两次：花体 $\g$ 表示李代数，斜体 $g$ 表示非阿贝尔耦合常数；$T$ 也出现两次：$\mathbb{T}$ 是标准三角形、$T$ 是欧几里得三角形（第 10 章还用 $T$ 表示三角剖分中的三角形）。译文保留原文字体区分。]
